\chapter{Implementation of the LC-KSVD Approach}\label{cha:chapter2}


\section{Mathematical Framework}
The LC-KSVD algorithm solves an optimization problem that balances reconstruction accuracy with label consistency and
classification error. Given a set of image patches $P$, the goal is to learn a dictionary $D$ and sparse coefficients
$X$ that minimize the following objective:

$$\min_{D, X, A, W} \|P - DX\|^2_F + \alpha\|Q - AX\|^2_F + \beta\|H - WX\|^2_F \quad \text{s.t.} \quad \forall i, \|x_i\|_0 \leq T$$

The first term $\|P - DX\|^2_F$ is the standard K-SVD reconstruction error. The second term $\|Q - AX\|^2_F$ introduces
the "discriminative sparse-code error," where $Q$ is a label-consistent constraint matrix that encourages patches from
the same class to activate the same dictionary atoms. The third term $\|H - WX\|^2_F$ is the classification error, where
$W$ is a linear classifier matrix.

This formulation ensures that the dictionary atoms are not just representative of the image data but are also class-specific.
In 3D CT multi-abnormality segmentation, this allows the model to learn distinct "sub-dictionaries" for each abnormality,
facilitating easier classification of individual voxels.


\section{Back-Projecting Class-Discriminative Atoms}
The unique intuition we propose in our approach is to generate spatial contribution maps through back-projection. Because
each atom in the dictionary $D$ is associated with a specific class label through the matrix $Q$, the model can theoretically
"explain" its decision for a test patch by examining which atoms were used in its sparse representation.

To generate a spatial contribution map for a class $c$, the algorithm identifies all atoms $d_k$ that have been assigned
to that class. For a query voxel or patch with sparse coefficients $x$, the contribution $C_c$ is calculated as the
partial reconstruction:

$$C_c = \sum_{k \in \text{Class } c} d_k x_k$$

This process should theoretically back-project the class-specific features into the original voxel space, creating a map
that highlights the regions contributing to the detection of a specific abnormality. for class-discriminative atoms
operating on raw voxel-space inputs, the back-projection is an exact, closed-form spatial attribution.


\section{Deriving Explainability through Segmentation Masks}
To evaluate the clinical accuracy of LC-KSVD, these contribution maps are converted into binary segmentation masks for
comparison with ground truth annotations. This conversion can be achieved through various thresholding methods.

Grad-CAM's approximation has two sources of imprecision — the gradient itself is a local linear approximation of a non-linear
function, and the spatial upsampling from the coarse feature map to input resolution introduces further interpolation error.
LC-KSVD's contribution map has only one source of imprecision: the thresholding step you apply to binarize it into a mask.
The map itself — before thresholding — is exact given the sparse code and the dictionary.


Unlike the passive Grad-CAM approach, atom back-projection is generative. The atoms are what the model actually used to
reconstruct and classify the signal. There's no proxy involved. The localization evidence is produced by the same
computation that produced the classification, not by a separate post-hoc analysis run afterwards.

%%%%%%%%%%%%%%%%%%%%%%%%%%%%%%%%%%%%%%%%%%%%%%%%%%%%%%%%%%%%%%%%%%%%%%%%%%%%%%%%%%%%%%%%%%%%%%%%%%
%%%%%%%%%%% END OF FILE
%%%%%%%%%%%%%%%%%%%%%%%%%%%%%%%%%%%%%%%%%%%%%%%%%%%%%%%%%%%%%%%%%%%%%%%%%%%%%%%%%%%%%%%%%%%%%%%%%%
